\section{Finite windows when a convergence modulus is unavailable}
\label{sec:metastability}

Some goals need a uniform tail bound: all later terms must be close, or
all later denominators must fail a Diophantine approximation test.  Other
goals use only one finite window, whose required length depends on its
starting point.  The latter demand is strictly weaker computationally.
We give a reusable theorem for it, including the case of represented
real values where exact comparison of the values is unavailable.

\subsection{Finite change budgets}

\begin{lemma}[A finite rise alternative]\label{lem:rise}
For a natural-valued sequence \(r_0,\ldots,r_n\), either
\[
 \exists i<n\quad r_{i+1}\le r_i,
\]
or \(r_n\ge r_0+n\).  In particular, if \(r_{B+1}\le B\), there is
an \(i\le B\) with \(r_{i+1}\le r_i\).
\end{lemma}

\begin{proof}
Keep a flag recording whether a nonrise has occurred in the scanned
prefix, its first index when present, and the current lower bound
\(r_i\ge r_0+i\) when absent.  At the next position compare the two
natural numbers.  A nonrise sets the flag; otherwise the new entry is at
least one larger.  This proves the alternative by bounded induction.
If there are \(B+1\) strict rises then
\(r_{B+1}\ge r_0+B+1>B\), contradicting the bound.
The flagged invariant is quantifier-free; the bounded existential in
the displayed conclusion is recovered from its stored index.
\end{proof}

Fix a PR control function \(g:\mathbb N\to\mathbb N\), and define
candidate starts
\begin{equation}\label{eq:window-starts}
 n_0=0,\qquad n_{i+1}=n_i+g(n_i)+1.
\end{equation}
The extra one ensures progress even if \(g(n_i)=0\).  The desired
window \([n_i,n_i+g(n_i)]\) lies within
\([n_i,n_{i+1}]\).

\begin{theorem}[Rational finite-window stability]\label{thm:rational-window}
Let \(s_n\) be a fixed PR nondecreasing rational sequence with
\(0\le s_n\le M\), for a supplied nonnegative rational \(M\).
For rational \(\varepsilon>0\), put
\(B=\lfloor M/\varepsilon\rfloor\).  Bounded search among
\(n_0,\ldots,n_B\) finds a start \(N\le n_B\) such that
\[
 N\le i,j\le N+g(N)
 \quad\Longrightarrow\quad |s_i-s_j|<\varepsilon.
\]
All functions and the proof are primitive recursive and available in
\ipra for the fixed sequence and control function.
\end{theorem}

\begin{proof}
The integers \(r_i=\lfloor s_{n_i}/\varepsilon\rfloor\) form a
nondecreasing sequence bounded by \(B\).  Lemma~\ref{lem:rise} gives
an index \(i\le B\) where two adjacent values coincide.  Monotonicity
places every intermediate value \(s_j/\varepsilon\) in the same
half-open interval \([r_i,r_i+1)\).  Therefore the difference between
any two such values is smaller than one.  Floors of rationals and the
finite search are PR.  The correctness proof uses finite-prefix
monotonicity and a Boolean record of the first successful comparison.
\end{proof}

The number of candidate windows depends only on the range and precision.
Their locations depend on the control \(g\).  Iterating a fixed PR
\(g\) for a variable natural number of steps is primitive recursion.
The statement is a schema for each such \(g\), or a theorem about a
finite supplied control table.  It is not a uniform PR evaluation
theorem for arbitrary codes of PR control programs.

\subsection{Soft endpoint selection for represented real sequences}

Taking the floor of an arbitrary represented real would reintroduce an
undecidable boundary test.  There is a different proof with explicit
slack and only rational comparisons.

\begin{theorem}[Represented finite-window stability]\label{thm:real-window}
Let \(s_n\) be a fixed PR family of regular real names.  Suppose arithmetic
proofs of monotonicity and \(0\le s_n\le B\) are supplied, where
\(B\) is a natural number.  Given \(k\), set
\[
 \varepsilon=2^{-k},\qquad
 K=2^{k+1}B+1,\qquad \delta=2^{-k-3}.
\]
Compute the candidate starts \eqref{eq:window-starts} through \(n_K\).
For each endpoint choose its prescribed stage-\(k+3\) rational value
\(a_i\), so \(|a_i-s_{n_i}|\le\delta\).  Let \(j<K\) be the first
index minimizing the rational difference \(a_{i+1}-a_i\).
Then \(N=n_j\) satisfies
\[
 N\le u,v\le N+g(N)
 \quad\Longrightarrow\quad |s_u-s_v|<2^{-k}.
\]
The search uses exactly \(K+1\) value observations, all at the same
precision, and at most \(K\) rational comparisons.  The conclusion has
the explicit margin
\( |s_u-s_v|\le2^{-k}-2^{-k}/(2K)\).
\end{theorem}

\begin{proof}
Write \(\Delta=a_{j+1}-a_j\).  Minimality and finite telescoping give
\[
 K\Delta\le\sum_{i<K}(a_{i+1}-a_i)
 =a_K-a_0\le B+2\delta.
\]
Consequently the true increase across the selected endpoints is at most
\[
 s_{n_{j+1}}-s_{n_j}
 \le\Delta+2\delta
 \le B/K+2\delta/K+2\delta
 \le B/K+4\delta<\varepsilon.
\]
The last strict inequality follows from
\(K>2B/\varepsilon\); it remains valid for \(B=0\), when \(K=1\).
Every value in the requested window lies between the selected endpoint
values by monotonicity, which proves the claim.

This proof never asks which of the real endpoint differences is smallest.
It minimizes their rational observations and retains enough error slack.
The telescoping and minimality assertions are finite rational formulas;
the monotonicity assumptions are instantiated only at the selected
endpoints and the two indices in the conclusion.  Monotonicity of
represented values remains a first-order hypothesis; bounding its index
range alone would not eliminate its unbounded precision quantifiers.
Each invocation can be expanded into the required finite-stage inequality
using the displayed positive margin and regularity.  In particular,
\(B/K+4\delta=\varepsilon-\varepsilon/(2K)\), so the strict
conclusion has a supplied rational separation.  The induction itself
concerns only the rational samples, their minimum, and their finite sum.
\end{proof}

This is often a useful replacement for a qualitative monotone-convergence
step when a subsequent proof asks for one stable finite block.  It is
not a proof that a primitive recursive convergence modulus exists.

\subsection{Cardinality information gives empty windows}

The Roth packet already explains the distinction between a cardinality
bound and a final denominator cutoff \cite{repoRoth}.  The following
general lemma makes its finite-window consequence a reusable contract.

\begin{theorem}[Bounded cardinality, finite empty window]\label{thm:empty-window}
Let \(G(a,n)\) be a fixed PR decidable predicate.  Suppose it has at most
\(B(a)\) distinct solutions, where \(B\) is a fixed PR function and a
proof of the cardinality assertion is supplied.
For a fixed PR function \(F(a,n)\ge n\), define
\[
 c_0=1,\qquad c_{i+1}=F(a,c_i)+1
 \quad(0\le i\le B(a)).
\]
Bounded searches find an \(i\le B(a)\) for which
\(G(a,n)\) fails at every \(n\in[c_i,c_{i+1})\).
In particular, the window contains \([c_i,F(a,c_i)]\).
\end{theorem}

\begin{proof}
The \(B(a)+1\) intervals are pairwise disjoint.  Scan them in order,
using a bounded search inside each.  If none is empty, store its first
solution.  The resulting list has \(B(a)+1\) distinct solutions,
contradicting the supplied cardinality theorem.  The search therefore
returns an empty interval.  The invariant records the selected solutions,
their validity, and the disjoint interval bounds; all are bounded tests.
The use of the cardinality premise is a final logical instantiation,
not induction on the universal assertion that there are few solutions.
\end{proof}

No general cardinality theorem for Roth or Faltings is proved by this
argument.  It identifies exactly what such a theorem would supply and
what it would not.

\subsection{Two explicit barriers to upgrading a window to a tail}

Let \(H(e,t)\) be the PR predicate asserting that machine \(e\) has
halted by step \(t\).  Define the single PR rational sequence
\begin{equation}\label{eq:specker}
 S_n=\sum_{e<n}4^{-(e+1)}H(e,n).
\end{equation}
It is nondecreasing: both its finite summation domain and each halting
indicator increase.  It is bounded above by the finite geometric bound
\(1/3\).  These facts use only bounded arithmetic.  Nevertheless it
has no computable Cauchy modulus.

Indeed, suppose a computable \(\mu(k)\) guaranteed
\(|S_m-S_n|\le2^{-k}\) for all \(m,n\ge\mu(k)\).
For a given machine \(e\), choose \(k=2e+3\) and
\(N=\max(\mu(k),e+1)\).  Simulate machine \(e\) through \(N\)
steps.  If it has not halted but halts later, the corresponding term in
\eqref{eq:specker} produces a later increase of at least
\(4^{-(e+1)}=2^{-2e-2}>2^{-k}\), contradicting the modulus.
Thus the finite simulation decides halting, which is impossible.
This is an external computability obstruction to a uniform construction,
not an assertion of an internal arithmetic proof of its own consistency.
The finite-window theorems still apply to \(S_n\).

Similarly, the predicate that \(t\) is the \emph{first} halting stage of
machine \(e\) has at most one solution.  A uniform computable cutoff
for all its solutions, obtained only from that cardinality bound, would
decide whether \(e\) ever halts by simulating up to the cutoff.  The
bound of one is therefore insufficient to produce a final cutoff, even
though it immediately yields the two-window instance of
Theorem~\ref{thm:empty-window}.

The practical lesson is precise.  When a downstream argument consumes
only a finite stable block, expose its control function and use the
finite-window contract.  When it consumes an eventual tail or a complete
solution list, preserve the missing modulus or height theorem as an
explicit remaining obligation.
